\section{Per-method protocols}\label{supp:protocols}

This section records five things for each method under test. It records the exact objects the
method consumes and produces, the reference it is scored against, its sampling or optimization
budget, the metric that turns its output into a faithfulness score, and the committed record
file (version-controlled, in the released repository) that holds the per-game results. All methods share one evaluation contract. We state that
contract once and then specialize it per method. Every method receives the same bit-exact VCS
state on each game of the scored set. We state the contract on the six core games
(\texttt{breakout}, \texttt{ms\_pacman}, \texttt{pong}, \texttt{qbert}, \texttt{seaquest},
\texttt{space\_invaders}), which carry the full tiered ground truth and the exact oracle and
serve as the running example; the aggregate leaderboard scores extend the identical protocol
to the 42-game scored set of \Sectionname\ref{supp:scoredset}. That state is not the boot screen.
It is a shared gameplay state, produced by a single seeded random-action stream (seed~$0$) run
for a $90$-frame prefix and then rolled forward for a $15$-frame analysis horizon. The stream is
deterministic, so every method sees the identical state on each game, and the run is asserted
bit-exact per game. An oracle cause-density gate screens the state before any method runs. The
gate counts how many candidate causes move the output above a floor and rejects states with too
few, so a near-flat oracle column cannot pass. This gate is what turns a degenerate frame into a
usable one; on \texttt{seaquest} it raises the count of live candidate causes from $2$ to $19$.
The candidate-cause set is the machine's own state. This set is the RAM bytes together with the
CPU registers and the TIA and RIOT registers exposed by the substrate. The output every method
attributes over is a shared screen-buffer region read at the end of the analysis horizon, not a
single hand-picked pixel. The reference is the intervention oracle of \Sectionname\ref{sec:oracle},
evaluated by exact re-runs of the emulator under $\mathrm{do}(u)$. A method that
cannot be applied to a given output regime is recorded as not-applicable rather than scored
as zero. An empty result is therefore never read as a wrong one. Faithfulness is the F
axis of the main-text triad (\Sectionname\ref{sec:triad}); we reference that definition and do not restate it.
Records live under \texttt{tools/xai\_study/<phase>/out/}; the aggregate they feed is
\texttt{tools/xai\_study/compare/out/leaderboard.json}.
Table~\ref{supp:leaderboard} consolidates the triad scores for all thirty methods and the
oracle positive control in one place; the per-method protocols follow.

\begin{table*}[t]
\centering
\caption{\textbf{Consolidated leaderboard: all methods, all triad axes.} One row per method
across the three phases, plus the intervention oracle as a positive control. $F$ is
faithfulness (Pearson correlation against the oracle, with the bootstrap-over-games 95\% CI in
brackets), $S$ is sufficiency, $M$ is minimality, all as defined in \Sectionname\ref{sec:triad};
Plaus is the elicited plausibility prior and $n$ is the number of scored games. Attribution
methods that emit only a saliency have no defined $S$ or $M$ and are marked ``---''; A5
(field potentials) and A7 (dimensionality reduction) likewise lack the $S$/$M$ decomposition.
Values are transcribed from \texttt{tools/xai\_study/compare/out/leaderboard.json} (the
committed numbers of record).}
\label{supp:leaderboard}
\footnotesize
\setlength{\tabcolsep}{4pt}
\begin{tabularx}{\textwidth}{@{}Xccccccc@{}}
\toprule
Method & Phase & Tradition & $F$ (95\% CI) & $S$ & $M$ & Plaus & $n$ \\
\midrule
A1 connectomics & A & intervention & 0.207 [0.130, 0.293] & 0.971 & 0.967 & 0.55 & 35 \\
A2 lesions & A & intervention & 0.965 [0.942, 0.983] & 0.591 & 1.000 & 0.55 & 42 \\
A3 tuning curves & A & correlational & 0.213 [0.136, 0.295] & 0.200 & 0.928 & 0.80 & 42 \\
A4 spike-word / pairwise & A & correlational & 0.225 [0.166, 0.283] & 0.133 & 0.363 & 0.80 & 34 \\
A5 field potentials & A & correlational & 0.376 [0.323, 0.428] & --- & --- & 0.80 & 40 \\
A6 Granger & A & correlational & 0.136 [0.091, 0.184] & 0.884 & 0.647 & 0.80 & 35 \\
A7 dim.\ reduction & A & dim.\ red. & 0.519 [0.457, 0.581] & --- & --- & 0.75 & 42 \\
A8 whole-state & A & descriptive & 1.000 [1.000, 1.000] & 1.000 & 0.075 & 0.30 & 42 \\
\midrule
Vanilla saliency & B & gradient & 0.380 [0.316, 0.445] & --- & --- & 0.90 & 42 \\
smoothgrad & B & gradient & 0.380 [0.317, 0.445] & --- & --- & 0.90 & 42 \\
Guided backprop & B & gradient & 0.322 [0.267, 0.375] & --- & --- & 0.90 & 42 \\
Grad$\times$Input / DeepLIFT & B & gradient & 0.486 [0.448, 0.526] & --- & --- & 0.90 & 42 \\
Integrated Gradients & B & gradient & 0.379 [0.314, 0.444] & --- & --- & 0.90 & 42 \\
Expected Gradients & B & gradient & 0.175 [0.110, 0.240] & --- & --- & 0.90 & 42 \\
kernelshap & B & gradient & 0.652 [0.596, 0.709] & --- & --- & 0.90 & 42 \\
lime & B & gradient & 0.644 [0.582, 0.705] & --- & --- & 0.90 & 42 \\
rise & B & gradient & 0.552 [0.500, 0.600] & --- & --- & 0.90 & 42 \\
occlusion & B & intervention & 0.687 [0.624, 0.746] & --- & --- & 0.55 & 42 \\
Extremal perturbation & B & intervention & 0.461 [0.410, 0.513] & --- & --- & 0.55 & 42 \\
On-distribution counterfactual & B & intervention & 0.433 [0.343, 0.522] & --- & --- & 0.55 & 42 \\
\midrule
Activation patching & C & causal & 1.000 [1.000, 1.000] & --- & --- & 0.50 & 42 \\
Causal scrubbing & C & causal & 0.979 [0.940, 1.000] & 0.979 & 0.814 & 0.50 & 42 \\
Interchange interventions (DAS) & C & causal & 1.000 [1.000, 1.000] & --- & --- & 0.50 & 42 \\
Logit / tuned lens & C & causal & 0.820 [0.768, 0.871] & --- & 0.224 & 0.50 & 42 \\
Path patching & C & causal & 0.580 [0.452, 0.710] & --- & --- & 0.50 & 42 \\
ACDC & C & causal & 0.625 [0.514, 0.730] & 0.299 & 1.000 & 0.50 & 35 \\
Attribution patching & C & gradient & 0.456 [0.345, 0.573] & --- & --- & 0.90 & 42 \\
Linear probing (+ control) & C & probing & 0.240 [0.153, 0.346] & --- & 0.932 & 0.85 & 42 \\
NMF / PCA dictionaries & C & dim.\ red. & 0.102 [0.049, 0.160] & 0.751 & 0.966 & 0.75 & 42 \\
Sparse autoencoder & C & dim.\ red. & 0.153 [0.096, 0.218] & 0.183 & 0.934 & 0.75 & 40 \\
\midrule
Intervention oracle & --- & oracle & 1.000 [1.000, 1.000] & 1.000 & 1.000 & 1.00 & 1 \\
\bottomrule
\end{tabularx}
\end{table*}

The neuroscience analyses (Phase~A) treat the running machine as a brain. It asks how far
each classical analysis recovers the known mechanism on the shared gameplay state. Each
method's baseline is the exact data-flow graph or causal role that the substrate makes
available. Its budget is the scored set (\Sectionname\ref{supp:scoredset}) unless noted,
with the six-game core as the running example, and its scoring metric is named in
Table~\ref{supp:tab:phaseA}. One core cell is honestly degenerate. On \texttt{ms\_pacman} the
pairwise analysis (A4) finds no coupled candidate pair to score at that frame, so its
\texttt{ms\_pacman} value is uninformative rather than a failure; the other five games carry
it.

\begin{table}[t]
\centering
\caption{\textbf{Phase~A protocols (neuroscience analyses).} One row per method. The columns
give the input it reads, the reference it is scored against, the sampling or intervention
budget, the scoring metric, and the committed record. All eight run on the scored set
(\Sectionname\ref{supp:scoredset}), with the six core games as the running example;
faithfulness is the F axis of \Sectionname\ref{sec:triad}. The whole-state recording (A8) is included as a
descriptive upper bound on coverage. Its faithfulness $F$ is trivial, so we read it for its
minimality $M$.}
\label{supp:tab:phaseA}
\scriptsize
\setlength{\tabcolsep}{4pt}
\begin{tabular}{@{}p{1.7cm}p{2.6cm}p{1.9cm}p{2.6cm}p{1.0cm}@{}}
\toprule
Method & Input / candidate causes & Reference & Scoring metric & Record \\
\midrule
A1 connectomics & RAM/register activity over trajectory & data-flow graph & data-flow graph $F_1$ vs.\ oracle & A1 \\
A2 lesions & single-cell knockouts via $\mathrm{do}(u)$ & true causal role & lesion-importance $\rho$ vs.\ role & A2 \\
A3 tuning curves & per-cell response vs.\ variable & variable coupling & $1-$ spurious-tuning rate & A3 \\
A4 spike-word / pairwise & pairwise cell correlations & true coupling & corr.\ $\rho$ vs.\ coupling & A4 \\
A5 field potentials & pooled-activity spectra & true causal share & $1-$ clock-explained var. & A5 \\
A6 Granger & directed lag structure & true data-flow & $1-$ false-edge rate & A6 \\
A7 dim.\ reduction & NMF/PCA over activity & known variables & NMF matched-component frac. & A7 \\
A8 whole-state & full state record & oracle decomp. & minimality $M$ vs.\ oracle & A8 \\
\bottomrule
\end{tabular}

\smallskip
\noindent\scriptsize Records: \texttt{tools/xai\_study/phaseA\_kording/out/A\{1..8\}\_*}; all
eight run on the 42-game scored set (\Sectionname\ref{supp:scoredset}), with the six core
games as the running example (\jkey{n_games} per method in \jkey{leaderboard.json}).
\end{table}

The attribution methods (Phase~B) run the everyday XAI toolbox over the same gameplay state.
These methods produce a real-valued saliency or attribution over the candidate-cause set.
We score that output by Pearson correlation against the oracle attribution. The gradient family
runs with the bilinear sub-pixel sampler turned on, so it is scored on a real position gradient
and not on the naive one. The naive gradient of the discrete position output is provably zero
(Prop.~\ref{prop:zero}); the sampler restores a finite gradient, but the restored gradient is
low and mixed in faithfulness rather than a fix. Expected Gradients draws its reference pool
from the states of the same seeded gameplay trajectory, not from a synthetic baseline. Two cells
are honestly degenerate and recorded as such. The on-distribution counterfactual returns a zero
delta on \texttt{space\_invaders}, \texttt{ms\_pacman}, and \texttt{qbert}, where its perturbation
does not move the shared output at that frame; those cells are marked not-valid rather than
scored. Expected Gradients scores low on the content regime of every game except \texttt{pong},
because the causal byte it should land on is static there. The budgets that matter are the
Monte-Carlo or perturbation counts of the sampling estimators, recorded per method in
Table~\ref{supp:tab:phaseB}. The seed-dispersion of the stochastic members is carried in
\texttt{leaderboard\_ci.csv} (\texttt{seed\_sd}, \texttt{seed\_sd\_source}) and reproduced
in \Sectionname\ref{supp:provenance}.

\begin{table}[t]
\centering
\caption{\textbf{Phase~B protocols (attribution methods).} Each method emits a saliency over
the candidate-cause set. We score it by Pearson correlation against the oracle attribution
(\Sectionname\ref{sec:oracle}). The reference is shared and the budget is the scored set of
\Sectionname\ref{supp:scoredset} (six core games as the running example). The budget column records the
sampling or perturbation count that drives each estimator. The seed column flags which
methods carry an independent-seed dispersion in \texttt{leaderboard\_ci.csv}.}
\label{supp:tab:phaseB}
\scriptsize
\setlength{\tabcolsep}{4pt}
\begin{tabular}{@{}p{2.9cm}p{1.75cm}p{4.3cm}p{1.15cm}@{}}
\toprule
Method & Tradition & Budget / hyperparameters & Seed sd \\
\midrule
Vanilla saliency & gradient & single backward pass & --- \\
smoothgrad & gradient & noise-averaged gradient; $\sigma$-robustness committed & 0.0 \\
Guided backprop & gradient & single guided pass & --- \\
Grad$\times$Input / DeepLIFT & gradient & DeepLIFT reference-difference & --- \\
Integrated Gradients & gradient & path integral, baseline = zero state & --- \\
Expected Gradients & gradient & expectation over baselines (refits) & yes \\
kernelshap & gradient & Monte-Carlo coalition sampling ($\propto 1/\sqrt{N}$) & 0.017 \\
lime & gradient & local linear fit, perturbation refits & 0.029 \\
rise & gradient & random-mask Monte-Carlo ($\propto 1/\sqrt{N}$) & 0.051 \\
occlusion & intervention & exhaustive single-cell masking & --- \\
Extremal perturbation & intervention & optimized mask, area-constrained & --- \\
On-distribution counterfactual & intervention & on-manifold $\mathrm{do}(u)$ resampling & --- \\
\bottomrule
\end{tabular}

\smallskip
\noindent\scriptsize Records: \texttt{tools/xai\_study/phaseB\_attribution/out/}; all twelve run
on the 42-game scored set (\Sectionname\ref{supp:scoredset}), six core games as the running
example, and are scored by Pearson correlation against the oracle (\Sectionname\ref{sec:oracle}). Seed
sd is the independent-seed dispersion from \texttt{leaderboard\_ci.csv}.
\end{table}

The mechanistic methods (Phase~C) are built to recover circuits. Each is scored
against the true data-flow or the exact recovered-minus-exact deviation. The causal members
of this phase are the ones that reach the oracle ceiling. We list their exact scoring metric
so that the ceiling is not read as a tautology. Each member is scored against the substrate's
known routine, and not against another method. Two members (NMF/PCA dictionaries and the
sparse autoencoder) are dimensionality-reduction methods carried here for comparison. The
sparse autoencoder has two aggregations that differ, which we reconcile in
\Sectionname\ref{supp:provenance}.

\begin{table}[t]
\centering
\caption{\textbf{Phase~C protocols (mechanistic methods).} Each method is scored against the
substrate's known circuit or the exact recovered-minus-exact deviation (\Sectionname\ref{sec:oracle}), and not
against another method. This is what keeps the oracle ceiling non-circular. The probing and
dictionary-learning members are included for comparison with the causal members.}
\label{supp:tab:phaseC}
\scriptsize
\setlength{\tabcolsep}{4pt}
\begin{tabular}{@{}p{3.4cm}p{1.4cm}p{5.2cm}@{}}
\toprule
Method & Tradition & Scoring metric \\
\midrule
Activation patching & causal & $\max|\text{recovered}-\text{exact}|$ vs.\ oracle \\
Causal scrubbing & causal & scrubbing-preserved performance \\
Interchange interventions (DAS) & causal & aligned interchange accuracy \\
Logit / tuned lens & causal & logit-lens fidelity to true intermediate \\
Path patching & causal & path-circuit $F_1$ vs.\ true routine \\
ACDC & causal & discovered-circuit best $F_1$ vs.\ data-flow \\
Attribution patching & gradient & corr.\ of approximation vs.\ exact patch \\
Linear probing (+ control) & probing & selectivity vs.\ control task \\
NMF / PCA dictionaries & dim.\ red. & matched-component fraction vs.\ vars \\
Sparse autoencoder & dim.\ red. & feature-variable matched fraction \\
\bottomrule
\end{tabular}

\smallskip
\noindent\scriptsize Records: \texttt{tools/xai\_study/phaseC\_mechanistic/out/}; all ten run
on the 42-game scored set (\Sectionname\ref{supp:scoredset}), with the six core games as the
running example.
\end{table}

The intervention oracle is the positive control, not a method under test. It is the exact
$\mathrm{do}(u)$ ground truth of \Sectionname\ref{sec:oracle}, evaluated by re-running the bit-exact machine. It
attains $F=S=M=1.0$ by construction. Its record triple is the intervention oracle, the
content-path gradient oracle, and their cross-check
(\texttt{tools/xai\_study/ground\_truth/out/}). It anchors the ceiling against which every
other row is read, which is why it occupies the thirty-first leaderboard row.
